\section{Background and Relation to Existing Standards}
\label{sec:background}

The representation of scientific and clinical evidence has been developed
across clinical informatics, biomedical ontology, and genomic knowledge
representation. Seven families of standards bear directly on our work: HL7 FHIR
Evidence~\cite{fhir_evidence}; the Evidence and Conclusion Ontology
(ECO)~\cite{eco2019} and the Scientific Evidence and Provenance Information
Ontology (SEPIO)~\cite{sepio}; the GA4GH Variant Annotation (VA) and Variation
Representation (VRS) specifications~\cite{ga4gh_va,ga4gh_vrs}; the ACMG/AMP
variant-interpretation guidelines~\cite{acmg2015}; the upper ontologies IAO
and OBI~\cite{iao2015,obi2016}; the PROV-O provenance model~\cite{prov_o};
and the ClinGen, MONDO, and HPO
curation and terminology resources~\cite{clingen2018,mondo,hpo}. Each captures
part of what a basic-science genetic publication reports, but none captures the
whole: a single conditional-knockout study may report molecular, cellular, and
clinical phenotypes at once, and a polygenic-score study may aggregate millions
of variants into one score. A per-family review, with the specific coverage
gaps that motivate our model, is given in Supplementary Note~\ref*{supp:sec:supp-background}. Of these
families, FHIR Evidence is the closest structural relative of our model, and we
position the model relative to it here.

\subsection{Relation to FHIR Evidence}
\label{sec:background-fhir}

FHIR is an HL7 information model rather than an OBO Foundry ontology. A
genetic-evidence model could, in principle, be built directly on the OBO
side, but doing so first would be premature. ECO and SEPIO provide,
respectively, an evidence-type vocabulary and a model of assertion
provenance and qualifiers. Neither provides the variable-and-certainty
content model for an evidence unit that FHIR Evidence provides and that
basic-science evidence needs to support translation from bench to bedside.

Aligning a genetics-only model directly to OBO before such a general content
model exists would yield a narrow resource. It would neither compose cleanly
with other basic-science domains nor reuse OBO's orthogonal vocabularies,
creating exactly the redundancy that OBO's emphasis on orthogonality and
reuse is intended to prevent.

We therefore take the structural content model from FHIR, generalize it from
clinical trials to basic science, and design the dimensions to align with
ECO, SEPIO, IAO, and OBI. Supplementary Note~\ref*{supp:sec:supp-crosswalk} provides a term-level
crosswalk to these vocabularies, including verified identifiers and the
remaining gaps. Genetics is the first specialization; extending the model to
other basic-science domains and completing the term-level alignment are
future work.

The HL7 FHIR Evidence, EvidenceVariable, and EvidenceReport
resources~\cite{fhir_evidence} were designed with systematic-review and
guideline-development workflows in mind, where evidence concerns an
intervention tested in a population. FHIR Evidence therefore has a different native
framing and does not by itself represent the evidence produced by basic and
pre-clinical research. Our three core classes address this gap.

The core classes are general across the basic sciences rather than specific
to genetics. They retain FHIR's separation between the evidence variable and
the certainty of a finding, incorporate FHIR's separate \texttt{statistic}
into \EvidenceVariable{} as the measured value, and replace the
clinical-trial content frame of population, intervention, comparator, and
outcome with a genetics-specific dimensional frame. Genetics is the
specialization developed in this paper (\cref{sec:model}):
\ScientificEvidence{}, \EvidenceVariable{}, and \EvidenceAssertion{}
specialize to \GeneticEvidence{}, \GeneticEvidenceVariable{}, and
\GeneticEvidenceAssertion{}, respectively. \Cref{tab:fhir-alignment} gives
the class-level correspondence. The field-level alignment, including the
certainty subcomponents discussed in \cref{sec:model-dimensions}, is given in
Supplementary Table~\ref*{supp:tab:fhir-fields}.
\footnote{We distinguish two notions that the literature often conflates.
The \texttt{assertion} property of a \ScientificEvidence{} item records the
object-level claim made by a paper. It is what the pilot records, and its
closest analogues are the GA4GH VA \texttt{Statement} and the SEPIO
\texttt{Assertion}, not the FHIR \texttt{Evidence.assertion} element, which
is a human-readable summary. \EvidenceAssertion{} is instead a
meta-predicate over decision rules that would use such
claims~\cite{our_preprint}. It has no direct analogue in FHIR, VA, or SEPIO and
is not instantiated in this paper. The property and the class share a name;
a future schema revision will resolve this collision.}

\begin{table}[tbp]
    \centering
    \caption{Class-level correspondence between the core model and FHIR Evidence
    R5. The alignment is architectural: the genetics-specific dimensional frame
    replaces FHIR's clinical-trial frame. Field-level detail is provided in
    Supplementary Table~\ref*{supp:tab:fhir-fields}.}
    \label{tab:fhir-alignment}
    \small
    \begin{tabularx}{\linewidth}{@{}llX@{}}
\toprule
This work & FHIR Evidence R5 & Relationship \\
\midrule
\ScientificEvidence
& \texttt{Evidence}
& Structural sibling; GEM replaces the PICO frame with a
  genetics-specific dimensional frame \\
\EvidenceVariable
& \texttt{EvidenceVariable} + \texttt{Evidence.statistic}
& GEM combines the variable definition and measured value in one class \\
\EvidenceAssertion
& no equivalent
& Meta-predicate over decision rules; part of the reasoning layer
  described in~\cite{our_preprint}, which is outside the scope of this paper \\
Credibility
& \texttt{Evidence.certainty}
& Structural parallel; both provide an overall rating with separable,
  genetics-specific subcomponents \\
\bottomrule
    \end{tabularx}
\end{table}
